% Manuscript for Part 3 (epic #51): the seed stencils of docs/stiff-features.md,
% written up for arXiv. Scaffold from #52; each section is drafted by the
% sub-issue named in its TRACE block.
%
% SOURCE-OF-TRUTH RULE: docs/stiff-features.md (the notes) stays canonical.
% This manuscript QUOTES the notes; it never becomes a second source of truth.
% Every number in this file must trace to a notes section or to the results
% cache of #54, and the assembly pass (#61) re-checks all of them. A TRACE
% comment names the notes section and the #27 sub-issue each part quotes.
%
% Section map (structure decided in the #51 grooming, 2026-09-20):
%   §1  Introduction              notes §1.1, §1.5, §3 (via #53)            #55
%   §2  Setting                   notes §1.1; dissertation / JCP 2017        #55
%   §3  Seeds in one dimension    notes §1.2–1.5 (#28, #31)                  #56
%   §4  1-D results               notes §2 (#32)                             #57
%   §5  Seeds in two dimensions   notes §4.1–4.2, §5.3 operator half         #58
%                                 (#34, #38, #39)
%   §6  2-D results               notes §5.2–5.6 (#37, #39–#42)              #59
%   §7  Discussion                #27 comments, notes §5.3, §5.6.1–5.6.2     #60
%   Figures and tables #54 · Literature #53 · Assembly #61 · arXiv #62
%
% Build: tectonic main.tex   (see paper/README.md)
\documentclass[11pt]{amsart}

\usepackage{amsmath,amssymb,amsthm}
\usepackage{graphicx}
\usepackage{booktabs}
\usepackage[colorlinks=true, linkcolor=blue, citecolor=blue, urlcolor=blue]{hyperref}
\graphicspath{{figures/}}

\newtheorem{theorem}{Theorem}[section]
\newtheorem{lemma}[theorem]{Lemma}
\newtheorem{proposition}[theorem]{Proposition}
\newtheorem{corollary}[theorem]{Corollary}
\theoremstyle{definition}
\newtheorem{definition}[theorem]{Definition}
\theoremstyle{remark}
\newtheorem{remark}[theorem]{Remark}

% Recurring quantities. Only macros the current text uses are defined; the
% section drafts add their own as they need them.
\newcommand{\Lu}{L_u}
\newcommand{\Lf}{L_f}
\newcommand{\px}{\partial_x}
\newcommand{\pt}{\partial_t}
\newcommand{\py}{\partial_y}
% Run-in headings inside a subsection (amsart's \paragraph adds its own
% period and no space above).
\newcommand{\para}[1]{\par\smallskip\noindent\textit{#1}.\enspace}
% Scientific notation for quoted measurements: $\sci{6.2}{-8}$.
\newcommand{\sci}[2]{#1 \times 10^{#2}}
% Draft stubs, visible in the PDF until the owning sub-issue replaces them.
\newcommand{\stub}[1]{\par\noindent\textit{[Stub, #1.]}\par}

% Working title; final wording is #61's call.
\title[Seed stencils for sub-grid material edges]{Seed stencils: high-order
finite differences and RBF-FD through material edges too steep for the grid}
\author{Bradley P. Martin}
\email{1bradley.martin1@gmail.com}
% Fixed revision date, updated deliberately with content changes. \today would
% make the committed main.pdf non-reproducible from the committed source.
\date{September 20, 2026}

\begin{document}

% ---------------------------------------------------------------------------
% ABSTRACT. Stub from the notes' summaries (README of the notes, CLAUDE.md);
% #55 revises it, #61 fixes the wording. Every claim below is in the notes:
%   knee at h = delta, seeds fourth order at every n     notes §2
%   ECT well-posedness, jump limit                        notes §1.3–1.4
%   straight edge -> triangular ODE systems               notes §4.1–4.2
%   stable at the standard gamma; rule "seed when         notes §5.3, §5.4
%     delta <= h"; 20x / 67x at 19,600 nodes
%   curved seeds = flat seeds at their own floor          notes §5.6
% ---------------------------------------------------------------------------
\begin{abstract}
Between a material whose parameters vary smoothly on the scale of the grid
and one with a jump lies the case of an edge that is smooth but changes over
a distance $\delta$ far smaller than the grid spacing $h$. Standard
high-order finite differences sample the coefficients at the nodes and lose
their order there until $h \lesssim \delta$, and the interface-aware
stencils built for a jump do not apply, since there is no jump. We build
stencils on the same equispaced (1-D) or scattered (2-D) nodes whose basis
functions are not monomials but \emph{seeds}: the $t = 0$ profiles of
solutions of the wave equation that are polynomial in time, continued
through the edge by ordinary differential equations in the coordinate normal
to it. The seeds form an extended complete Chebyshev system for every
$\delta$, so the stencil solve is nonsingular, and as $\delta \to 0$ they
reduce to the translated polynomial basis of the interface-aware stencils,
so one construction covers the jump, the sub-grid edge and the resolved
edge. In one dimension, standard fourth-order finite differences have a
knee at $h = \delta$ while the seed stencils are fourth order at every
resolution with the errors of the jump case. In two dimensions, for the
elastic wave equation discretised with radial basis function-generated
finite differences (RBF-FD) on scattered nodes, the seeds of a straight
edge reduce to triangular systems of ordinary differential equations in the
normal coordinate; the seed-augmented stencils are stable at the standard
hyperviscosity, cut the error through an unresolved edge by up to an order
of magnitude or more on the node sets tested, and lead to the rule: seed
when $\delta \le h$. The same seeds marched along the true normal handle a
curved edge at the same accuracy on those node sets.
\end{abstract}

\maketitle

% ===========================================================================
\section{Introduction}
\label{sec:intro}
% TRACE: notes §1.1 (the setting, the "twilight zone"), §1.2 (the chain,
% condition numbers 20–60), §1.3 (the jump limit: 10% at delta = h/10, 1% at
% h/100 for a stencil of half-width 0.02), §1.4 (ECT, the truncation-error
% argument), §1.5 (the Taylor route, poles at ±i pi delta / 2), §2 (the knee:
% rates 1.3 and 2.2, the factor 200 between n = 400 and 800 at
% delta = 0.0025; seeds match the jump case to three digits), §4.1–4.2 (the
% straight-edge reduction; 30 seeds, 600 states), §5.4 (20x in v, 67x in
% spurious u at 19,600 nodes; crossover at h ≈ delta; 1.4–2.0x worse on a
% resolved edge; the rule), §5.5 (1.4–1.9x at 26.6°; the ablation worse than
% naive; rates 2.5, 2.3, 1.9), §5.6 (curved: jump stencils at the floor,
% seeds 2.2–7x and 3–12x below naive). Drafted by #55. The prior-work
% paragraph is #53's (LITERATURE.md): cite, don't claim.

High-order finite differences for wave propagation are in good shape at
both ends of the range of media. When the material parameters vary
smoothly on the scale of the grid, sampling them at the nodes costs
nothing and a fourth-order stencil stays fourth order. When they jump,
the stencils that straddle the jump can be rebuilt from piecewise
polynomials that satisfy the interface conditions, and the scheme keeps
its order on a grid that takes no notice of the interface
\cite{Martin2016,MartinFornberg2017}. Between the two lies a case that is
neither. An edge can be smooth, so that no interface condition applies,
and still change over a distance $\delta$ smaller than the grid spacing
$h$. Any smooth model contains such edges once it is discretised coarsely
enough, and any jump becomes one under the slightest smoothing.

A standard stencil through such an edge samples the coefficients at its
nodes, sees a jump, and treats it as if the coefficients were smooth in
between, which is the first-order failure of the jump case. It recovers
only once $h \lesssim \delta$. In the one-dimensional experiment of
Section~\ref{sec:results1d}, fourth-order finite differences (FD4)
through an edge of width $\delta = 0.0025$ converge at rates $1.3$ and
$2.2$ while $h > \delta$, drop by a factor $200$ between $n = 400$ and
$n = 800$ nodes ($h = 2\delta$ to $h = \delta$), and are fourth order
after that: a knee at $h = \delta$. The two existing repairs do not
reach this regime. The interface-aware stencils built for a jump are the
$\delta \to 0$ limit of the right construction, as
Section~\ref{sec:seeds1d} shows, but their weights are off at first order
in $\delta / h$, by $10\%$ of the largest weight at $\delta = h/10$ and
$1\%$ at $\delta = h/100$ in the example of Section~\ref{sec:jumplimit}.
The other route expands the coefficients in Taylor series about the
edge, as the dissertation's construction can for smoothly varying
coefficients \cite{Martin2016} and as polynomial quasi-Trefftz spaces do
in general \cite{ImbertGerardMoiolaStocker2023}; but a profile such as
$\tanh(x / \delta)$ has poles at $\pm i \pi \delta / 2$, so the expansion
converges only within $\pi \delta / 2$ of the edge while the stencil
reaches $2h \gg \delta$.

The construction of this paper keeps everything about the scheme except
the basis functions of the stencils that see the edge. In place of the
monomials $1, x, x^2, \ldots$ a stencil is made exact on \emph{seeds}:
the $t = 0$ profiles of solutions of the wave equation that are
polynomial in time. For the operator $\Lu = \rho^{-1} \px K \px$ that
governs the particle velocity, such a profile satisfies $\Lu^m g = 0$
for some $m$, and with constant coefficients these profiles are exactly
the polynomials. Anchored at the stencil's evaluation point $x_e$ they
form a chain, $\phi_0 = 1$, $K \phi_1' = K(x_e)$,
$\Lu \phi_k = k (k - 1) c_e^2 \phi_{k-2}$ with $\phi_k(x_e) = \phi_k'(x_e)
= 0$, which is integrated through the edge as a first-order system of
ordinary differential equations that never differentiates the
coefficients: one short march per stencil that sees the edge. Three
facts make this a method rather than a trick. The seeds form an extended
complete Chebyshev system for every $\delta$, so the per-stencil solve is
nonsingular on any distinct nodes, including at $\delta = 0$ where the
seeds have kinks, with the conditioning of a Vandermonde matrix
(condition numbers $20$ to $60$ from $\delta = 10^{-5}$ to $1$). The
truncation error of the seed stencil on a solution is
$O(h^4\, \px \pt^4 u)$, bounded independently of $\delta$, where FD4's
$O(h^4 u^{(5)})$ carries $u^{(5)} \sim \delta^{-4}$ inside the edge. And
the chain read in the weak sense across a jump is the translated Taylor
basis of the interface-aware stencils, so one construction covers the
jump (algebra), the sub-grid edge (an ODE march) and the resolved edge,
where the seeds and the monomials give different weights, first order in
$h / \delta$ apart, and both are fourth order on solutions.

The same idea carries to two dimensions, to the elastic wave equation
on scattered nodes discretised with radial basis function-generated
finite differences (RBF-FD). For a straight edge the material depends on
the normal coordinate only, and an ansatz polynomial in the tangential
coordinate turns the chain into triangular systems of one-dimensional
ODEs in the normal coordinate: all $30$ seeds of one stencil, the
monomials to degree four in each velocity component, march as one linear
system of $600$ states anchored at the evaluation node. The seeds
replace the polynomial augmentation of the interface-aware RBF-FD
stencils; the Gaussian part of the stencil and the coupled saddle-point
solve stay as they are. For a curved edge the same seeds are marched
along the true normal through the stencil's foot point, an approximation
of zeroth order in the curvature, which is the one the jump stencils
already make.

\subsection*{Contributions}
\begin{enumerate}
\item In one dimension: the seed chain, its nested kernels, its
  well-posedness as an extended complete Chebyshev system, and the jump
  construction as its $\delta \to 0$ limit (Section~\ref{sec:seeds1d}).
  On the same equispaced grid the seed stencils are fourth order at every
  resolution, with errors matching the jump case to three digits, while
  FD4 has its knee at $h = \delta$ (Section~\ref{sec:results1d}).
\item In two dimensions: the elastic seeds of a straight edge as
  triangular systems of ODEs in the normal coordinate, marched as one
  system per stencil; their assembly into seed rows for the elastic
  operator and hyperviscosity rows that annihilate the seeds; and the
  finding that the seeded operator is stable at the standard
  hyperviscosity at every edge width (Sections~\ref{sec:seeds2d}
  and~\ref{sec:results2d}).
\item A rule: \emph{seed when $\delta \le h$}. Through an edge the nodes
  never resolve the seeds are fourth order at every resolution tested,
  $20\times$ below the standard scheme in the velocity and $67\times$ in
  the spurious tangential velocity at $19{,}600$ nodes, at the seed
  operator's own resolution floor; the crossover is at $h \approx
  \delta$; and through a resolved edge the seed stencils are $1.4$ to
  $2\times$ worse than the standard ones, so a resolved edge should not
  be seeded.
\item At oblique incidence the tangential dependence of the seeds is
  what acts: the full seeds beat the standard scheme by $1.4$ to
  $1.9\times$ and the normal-only ablation is worse than doing nothing
  through a sharp edge. Every scheme converges at about order $2.5$
  there, the jump-aware stencils included, because the mode-converted
  shear waves sit at a pre-asymptotic floor on these node sets; the
  result is reported as it is.
\item Through a curved edge the jump stencils sit at the resolution
  floor at every node count, and the seeds marched along the true normal
  equal the flat seeds at their own floor, $2.2$ to $7\times$ below the
  standard scheme at $\delta = 0.005$ and $3$ to $12\times$ at
  $\delta = 0.0025$ on node sets up to $19{,}600$ nodes; the zeroth-order
  geometry first shows at the finest of them.
\item An open implementation in Python with a test for every numerical
  routine, in which every figure and table of this paper is regenerated
  by a committed script from a results cache.
\end{enumerate}

\subsection{Relation to prior work}
\label{sec:priorwork}
\stub{\#53 supplies this paragraph; \#55 places it}

\subsection*{Outline}
Section~\ref{sec:setting} fixes the equations, the discretisations and
the interface-aware stencils of the jump case that the seeds extend.
Section~\ref{sec:seeds1d} constructs the seeds in one dimension and
proves what can be proved; Section~\ref{sec:results1d} is the
one-dimensional experiment. Sections~\ref{sec:seeds2d}
and~\ref{sec:results2d} do the same in two dimensions, flat edges first,
then oblique incidence and curved edges. Section~\ref{sec:discussion}
records the limitations and what we would do next.

% ===========================================================================
\section{Setting}
\label{sec:setting}
% TRACE: notes §1.1 (the 1-D system, the layer, the tanh profile with
% periodic images, CFL 0.4), §5.1 (the 2-D blend: lambda, mu, rho linear in
% the tanh weight; the signed normal distance), §2 (errors: relative l2 in f
% at t = 1); code defaults: wave1d/domain.py (LayeredMedium: (c, rho) =
% (1, 1) / (2, 1), layer [0, 0.5)), wave1d/operators.py (order 4),
% wave2d/domain.py (LayeredMedium2D: (1, 1, 1) / (4, 4, 2), y0 = 0.25, 0.5;
% straddling_rows: 3 rows per side at (1/2 + j sqrt 3 / 2) h),
% wave2d/rbf.py (Gaussians, eps = 0.4 / d with d the third-nearest
% neighbour), wave2d/operators.py (30 / degree 4; 19 / degree 3 across;
% hyper_power 3; gamma = 2.4e-11 (h / 0.02)^5), wave2d/simulate.py (CFL 0.5
% with the hyperviscosity cap), scripts/wave2d_stiff.py (sharpness 15,
% centre 0.875, t = 1, rel_error, rates in h = n^{-1/2}); the dissertation
% §2.1 (eq. 1–3, 22–28) and §3.3, JCP 2017 pp. 23–33 (Type 3 stencils,
% eq. 42, 30/4 and 19/3, the curvature remark), quoted not re-derived;
% CLAUDE.md's port notes for the scope. Drafted by #55.
%
% NOTATION fixed here for the later sections (#56–#60): n nodes, h the
% spacing (2/n in 1-D, n^{-1/2} in 2-D), delta the edge width; 1-D fields
% (u, f), K = rho c^2, L_u and L_f; 2-D velocities (u, v) and stresses
% sigma_xx, sigma_xy, sigma_yy (the code's f, g, h; h is reserved for the
% spacing), K = lambda + 2 mu; the local frame (x', y') with y' normal;
% the anchor x_e; seeds phi_k in 1-D.

\subsection{One dimension}
\label{sec:setting1d}
The one-dimensional problem is the two-way wave equation in first-order
form on the periodic interval $[-1, 1)$, with $u$ the particle velocity
and $f$ the stress,
\begin{equation}
  \label{eq:sys1d}
  \rho\, u_t = f_x, \qquad f_t = K\, u_x, \qquad K = \rho c^2,
\end{equation}
where $c(x) > 0$ is the wave speed and $\rho(x) > 0$ the density. Eliminating one
field gives a second-order operator for each,
\begin{equation}
  \label{eq:Lops}
  u_{tt} = \Lu u, \quad \Lu = \frac{1}{\rho} \px K \px; \qquad
  f_{tt} = \Lf f, \quad \Lf = K \px \frac{1}{\rho} \px .
\end{equation}
The medium is a background material $(c, \rho) = (1, 1)$ with one layer
$[x_1, x_2) = [0, \tfrac12)$ of $(c, \rho) = (2, 1)$, so the contrast in
$K$ is $4$ and the impedance $\rho c$ doubles. With edge width
$\delta > 0$ the two edges of the layer are $\tanh$ transitions,
\begin{equation}
  \label{eq:profile1d}
  \begin{gathered}
  s(x) = \frac12 \sum_{m \in \mathbb{Z}} \Bigl[
    \tanh\frac{x - x_1 + 2m}{\delta} - \tanh\frac{x - x_2 + 2m}{\delta}
  \Bigr], \\
  c = c_{\mathrm{bg}} + (c_{\mathrm{layer}} - c_{\mathrm{bg}})\, s,
  \qquad \rho \text{ likewise},
  \end{gathered}
\end{equation}
summed over the periodic images so that the profile is smooth and
periodic to rounding; $\delta = 0$ is the jump. For
$\delta \ll h$ a stencil sees a jump between two of its nodes; for
$\delta \gg h$ it sees a smooth function; the case in between is the
subject of the paper.

The grid has $n$ cell-centred nodes, $x_i = -1 + (i + \tfrac12) h$ with
$h = 2 / n$, and the semi-discrete system is
\begin{equation}
  \label{eq:semi1d}
  \rho_i\, \dot u_i = (D_f f)_i, \qquad \dot f_i = K_i\, (D_u u)_i,
\end{equation}
with $\rho_i = \rho(x_i)$, $K_i = K(x_i)$ and $D_u$, $D_f$ two
differentiation matrices with the sparsity of the centred five-point
stencil. In the \emph{naive} scheme both are the standard fourth-order
matrix, Fornberg's weights on the monomials \cite{Fornberg1988}: the
coefficients are sampled at the nodes and nothing else is done about the
edge. In the \emph{seed} scheme the rows whose stencils see a varying
material are rebuilt as in Section~\ref{sec:seeds1d}, separately for
$u$ and for $f$ because the two fields carry different continuity
conditions; every other row is the standard one. Time stepping is the
classical fourth-order Runge--Kutta method (RK4) with
$\Delta t = 0.4\, h / c_{\max}$. Errors are relative $\ell^2$ errors in
the stress $f$ at $t = 1$ against a reference solution
(Section~\ref{sec:results1d}).

\subsection{Two dimensions}
\label{sec:setting2d}
The two-dimensional problem is the isotropic elastic wave equation in
first-order form on the doubly periodic unit square, with velocities
$(u, v)$ and stresses $(\sigma_{xx}, \sigma_{xy}, \sigma_{yy})$,
\begin{equation}
  \label{eq:sys2d}
  \begin{aligned}
    \rho\, u_t &= \px \sigma_{xx} + \py \sigma_{xy}, &
    \partial_t \sigma_{xx} &= (\lambda + 2\mu)\, u_x + \lambda\, v_y, \\
    \rho\, v_t &= \px \sigma_{xy} + \py \sigma_{yy}, &
    \partial_t \sigma_{xy} &= \mu\, (u_y + v_x), \\
    & & \partial_t \sigma_{yy} &= \lambda\, u_x + (\lambda + 2\mu)\, v_y,
  \end{aligned}
\end{equation}
with Lam\'e parameters $\lambda, \mu > 0$ and density $\rho > 0$, P- and
S-wave speeds $c_p = \sqrt{(\lambda + 2\mu) / \rho}$ and
$c_s = \sqrt{\mu / \rho}$, and $K = \lambda + 2\mu$ in the role of
$\rho c^2$. (The implementation and the dissertation name the stresses
$f, g, h$; we write $\sigma$ so that $h$ can stay the node spacing.)
The medium is a background $(\lambda, \mu, \rho) = (1, 1, 1)$ with a band
$(4, 4, 2)$ between the two interfaces $y = y_0 + A \sin 2\pi x$,
$y_0 = \tfrac14$ and $\tfrac12$, so $c_p$ rises from $\sqrt 3$ to
$\sqrt 6$ and $c_s$ from $1$ to $\sqrt 2$ across each interface. The
flat case is $A = 0$ and the curved case $A = 0.02$, the two geometries
of the dissertation's elastic tests \cite{Martin2016}. With edge width
$\delta > 0$ each interface becomes a $\tanh$ transition in the signed
normal distance $d$ to it, with the band's weight $s$ summed over the
periodic images in $y$,
\begin{equation}
  \label{eq:profile2d}
  \begin{gathered}
  s = \tfrac12 \bigl[ 1 + \tanh(d / \delta) \bigr], \\
  (\lambda, \mu, \rho) = (\lambda, \mu, \rho)_{\mathrm{bg}}
  + \bigl[ (\lambda, \mu, \rho)_{\mathrm{band}}
  - (\lambda, \mu, \rho)_{\mathrm{bg}} \bigr]\, s,
  \end{gathered}
\end{equation}
so that $\lambda$, $\mu$ and $\rho$, and hence $K$, are linear in the
weight while $c_p$ and $c_s$ are not. For a flat interface $d$ is the
vertical offset; for a curved one it is the true signed distance,
through the foot point on the curve. Two edges a gap $g$ apart reach
only $\tanh(g / 2\delta)$ of the contrast between them, so the band's
width of $\tfrac14$ bounds the edge widths we can pose; nothing here
approaches that bound.

\para{Node sets}
A node set of $n$ nodes is built as in \cite{MartinFornberg2017}: six
fixed rows of nodes hug each interface, three on each side at normal
distances $(\tfrac12 + j \tfrac{\sqrt 3}{2}) h$, $j = 0, 1, 2$, with
alternate rows shifted by half a spacing so that the rows straddle the
interface in a hexagonal pattern, and every other node is placed by a
simulated electrostatic repulsion from a random start. The straddling
rows were found empirically to be what keeps the interface-aware
stencils stable, and we keep them for the seed stencils. Throughout,
$h = n^{-1/2}$ denotes the mean spacing, and convergence rates are
measured in it; the node counts are $n = 2500$, $4900$, $10{,}000$ and
$19{,}600$, so $h$ runs from $0.02$ to $0.0071$.

\para{RBF-FD}
For a stencil of $s$ nodes $x_1, \ldots, x_s$ about an evaluation node
$x_c$, the weights $w$ of a linear operator $\mathcal{L}$ solve the
saddle-point system
\begin{equation}
  \label{eq:rbffd}
  \begin{bmatrix} A & P \\ P^{\mathsf T} & 0 \end{bmatrix}
  \begin{bmatrix} w \\ w_* \end{bmatrix}
  =
  \begin{bmatrix}
    \mathcal{L}\, \varphi(|x - x_i|) \big|_{x_c} \\
    \mathcal{L}\, p_m(x) \big|_{x_c}
  \end{bmatrix},
\end{equation}
with $A_{ij} = \varphi(|x_i - x_j|)$ for the Gaussian
$\varphi(r) = e^{-(\varepsilon r)^2}$, $P_{im} = p_m(x_i)$ over the
bivariate monomials $p_m$ up to a chosen degree, evaluated in
coordinates scaled by the stencil radius, and $w_*$ discarded. The constraint rows make the weights
exact on the polynomials, which take over from the Gaussians under
refinement, and the shape parameter is scaled to the local spacing,
$\varepsilon = 0.4 / d$ with $d$ the distance from $x_c$ to its
third-nearest neighbour.\footnote{The published description
\cite{MartinFornberg2017} says the nearest neighbour; the original
\textsc{Matlab} and the present implementation use the third-nearest.
The choice fixes a constant in $\varepsilon h$ and nothing else.} Away
from the interfaces every stencil has $30$ nodes and is exact on the
polynomials of degree $\le 4$; stencils across a jump have $19$ nodes
and degree $\le 3$, with the piecewise basis of
Section~\ref{sec:jumpstencils}. Two matrices $D_x$ and $D_y$ result, and
the $5n \times 5n$ block operator of \eqref{eq:sys2d} is assembled from
them with the material coefficients evaluated at each row's node. A
small hyperviscosity $\gamma \Delta^3$ is added to every field to move
the spurious eigenvalues of the RBF-FD operator into the left half-plane
\cite{FornbergLehto2011}: the $\Delta^3$ stencils are built on the same
nodes and the same basis as the derivative stencils, and
$\gamma = \gamma_0\, (h / h_0)^5$ with $\gamma_0 = 2.4 \times 10^{-11}$
at $h_0 = 0.02$, which keeps $\gamma \Delta^3$ proportional to $1 / h$
like the wave operator so that one CFL number serves every $n$. Time
stepping is RK4 with $\Delta t = 0.5\, h / c_{p,\max}$, capped where
necessary by the hyperviscosity spectrum.

\para{Pulses and errors}
The normal-incidence runs send a plane P-wave pulse, Gaussian in $y$
with sharpness $15$, from $y = 0.875$ in the background towards $-y$
through both interfaces; the true solution then has $u \equiv 0$, and
the largest $|u|$ a scheme produces is a clean measure of what its
stencils do at the edges. The oblique runs send a plane-wave train at
$26.6^\circ$ to the normal (direction $(1, 2)$), which mode-converts at
each interface. Errors are relative $\ell^2$ errors in $v$ at $t = 1$
against the references described with each experiment in
Section~\ref{sec:results2d}, together with the spurious $u$ at normal
incidence and the relative error in $u$ otherwise.

\subsection{Interface-aware stencils for a jump}
\label{sec:jumpstencils}
The seeds extend the interface-aware stencils of
\cite{Martin2016,MartinFornberg2017}, which we summarise in the form the
later sections use; the derivations are in those references.

\para{One dimension}
Write \eqref{eq:sys1d} as $\pt (u, f)^{\mathsf T} = \mathcal{D}\, (u,
f)^{\mathsf T}$ with the block operator
\[
  \mathcal{D} = \begin{bmatrix} 0 & \rho^{-1} \px \\ K \px & 0 \end{bmatrix}.
\]
Across a material jump, $u$ and $f$ are continuous for all time,
hence so are all their time derivatives, and $\pt^k (u, f)^{\mathsf T} =
\mathcal{D}^k (u, f)^{\mathsf T}$ converts each of those conditions into
a condition on spatial derivatives, with the material of the side it is
evaluated on: $\mathcal{D}_L^k (u_L, f_L)^{\mathsf T} = \mathcal{D}_R^k
(u_R, f_R)^{\mathsf T}$ at the jump for $k = 0, 1, 2, \ldots$ Expanding
$u$ and $f$ on each side as polynomials of degree four about the jump
and imposing these conditions for $k \le 4$ gives a square
\emph{continuity matrix} per side, $C_L a_L = C_R a_R$ on the
coefficient vectors. Keeping the monomials on one side and translating
them across with $C_R^{-1} C_L$ yields the \emph{translated basis}:
piecewise polynomials, one per monomial, that satisfy every continuity
condition. A stencil whose nodes straddle both edges of a layer thinner
than the stencil chains two translations. The stencil weights are then
the ones that differentiate every basis function exactly at the
evaluation point, a small dense solve per stencil, and the $f$ basis is
built by the same construction rather than from the $u$ basis through
the operator, which would lose an order. This is the scheme that
Section~\ref{sec:results1d} calls the jump case; the even continuity
conditions relate a field to itself and the odd ones route through the
other field, a pattern the seeds inherit.

\para{Two dimensions}
A stencil near an interface works in a local frame: origin at the
nearest interface point, $x'$ along the tangent and $y'$ along the
normal. The isotropic system \eqref{eq:sys2d} keeps its form under
rotation, so in that frame the rotated fields obey the same equations
and the interface is, locally, $y' = 0$. The velocities $(u', v')$ and
the tractions $(\sigma'_{x'y'}, \sigma'_{y'y'})$ are continuous across
it for all time; the remaining stress $\sigma'_{x'x'}$ may jump. Expand
$(u', v')$ on each side in the monomials of $(x', y')$ up to degree
four, convert the time derivatives of the continuous quantities through
powers of the block operator (even powers relate velocities to
velocities, odd powers route through the stresses), and match the
coefficients of the powers of $x'$ at each order: a square continuity
matrix per side, a translated velocity basis, truncated to degree three.
The stress basis is generated from it by applying the stress rows of the
operator on each side and dropping the three columns that come from
rigid motions, $27$ functions in all. The RBF-FD solve
\eqref{eq:rbffd} then couples the fields: the stress rates need weights
on $u'$ and $v'$ data jointly, the velocity rates weights on the three
stresses jointly, with the Gaussian block repeated on the diagonal and
the piecewise basis in place of $P$. The weights are rotated back to the
laboratory frame, and the hyperviscosity stencils across the interface
use the same piecewise basis, which the original work found necessary
for accuracy. The interface is treated as locally flat in the frame;
carrying curvature into the continuity conditions by expanding the
interface in $x'$ is possible \cite{MartinFornberg2017} but was not used
by the original code and is not used here. This is the approximation we
call zeroth order in the curvature, and the curved seeds of
Section~\ref{sec:seeds2d} make the same one.

\subsection{Scope of the implementation}
\label{sec:scope}
The experiments run on a Python reimplementation of the 2016--2017
methods (NumPy and SciPy, with a test for every numerical routine) that
reproduces the constructions above and restructures the code. Of the
original two-dimensional code it leaves out three things: the
\emph{triple stencils} that see both interfaces of a band thinner than
one stencil (the band here is $\tfrac14$ wide and none is needed), the
variable Lam\'e parameters inside the band of the published test case
(the band here is a constant material), and the finite-difference far
field of the FD/RBF-FD hybrid (every stencil here is RBF-FD). Point
sources and the Marmousi model are out of scope. In one dimension the
implementation takes piecewise-constant materials, with the smooth edges
of \eqref{eq:profile1d}, and not the dissertation's smoothly varying
background, because the exact reference of the jump case, a ray sum,
needs constant pieces. Everything is fourth order (FD4 in one dimension,
degree four away from and degree three across the edges in two) at one
material contrast; Section~\ref{sec:discussion} returns to both
restrictions.
% ===========================================================================
\section{Seeds in one dimension}
\label{sec:seeds1d}
% TRACE: notes §1.2 (seeds as profiles of time-polynomial solutions, the
% chain), §1.3 (nested kernels, the jump limit), §1.4 (ECT well-posedness),
% §1.5 (the Taylor route); code wave1d/stiff.py (#28, #31). Drafted by #56.
% Figure: the 1-D seed basis (#54; today docs/figures/wave1d_stiff_seeds.png).
\subsection{Seeds as the profiles of solutions polynomial in time}
\stub{\#56: $u_{tt} = \Lu u$ with $\Lu = \rho^{-1} \px K \px$; the chain
$\phi_0 = 1$, $K \phi_1' = K(x_e)$, $\Lu \phi_k = k(k-1) c_e^2 \phi_{k-2}$
with $\phi_k(x_e) = \phi_k'(x_e) = 0$; the even/odd ($u$/$f$) alternation;
first-order form so that $K$ is never differentiated}
\subsection{Nested kernels and the jump as the limit}
\label{sec:jumplimit}
\stub{\#56: $\{1\} \subset \ker \Lu \subset \{\Lu g = \mathrm{const}\}
\subset \ker \Lu^2 \subset \cdots$; the seed weights converge to the
interface-aware weights at first order in $\delta$}
\subsection{Well-posedness}
\stub{\#56: the seeds are an extended complete Chebyshev system (P\'olya,
Karlin--Studden), so the per-stencil solve is nonsingular on any distinct
nodes for every $\delta$}
\subsection{Why the Taylor route does not reach a sub-grid edge}
\stub{\#56: notes §1.5; the quasi-Trefftz comparison}
\subsection{Stencil assembly and cost}
\stub{\#56: the Fornberg-style solve on the seeds; one ODE march per
stencil}

% ===========================================================================
\section{One-dimensional results}
\label{sec:results1d}
% TRACE: notes §2 (#32), scripts/wave1d_stiff.py; numbers from #54's cache.
% Drafted by #57. Figures: knee (convergence), snapshot; table: the knee.
\stub{\#57: test problem (medium, contrast, pulse, $\delta \in \{0, 0.0025,
0.01\}$); the Fourier pseudo-spectral reference; the knee of standard FD4 at
$h = \delta$; seeds fourth order at every $n$ with the jump case's errors;
condition numbers $20$--$60$ from $\delta = 10^{-5}$ to $1$}

% ===========================================================================
\section{Seeds in two dimensions}
\label{sec:seeds2d}
% TRACE: notes §4.1 (design, #34), §4.2 (the elastic seeds of a straight
% feature, #38), the operator half of §5.3 (#39), §5.6 for route (a);
% code wave2d/seeds.py, wave2d/operators.py. Drafted by #58. Figure: 2-D seed
% cross-sections (#54, new).
\subsection{Design: a straight edge}
\stub{\#58: the $x'$-polynomial ansatz turns the 2-D chain into triangular
systems of 1-D ODEs in the normal coordinate; tangential dependence stays
polynomial; tractions as flux variables}
\subsection{The elastic seeds}
\stub{\#58: all 30 seeds of one stencil marched as one 600-state system
anchored at the evaluation node; constant-material, jump-limit, residual and
conditioning checks}
\subsection{Stencil assembly}
\stub{\#58: 19-node seed rows for the elastic operator; $\Delta^3$ rows on
the naive 30-node footprint with the seeds annihilated; the seeded region
($19\delta$) and its trim}
\subsection{Curved edges}
\stub{\#58: route (a), the seeds marched along the true normal through the
foot point, zeroth order in curvature, the same approximation as the jump
stencils; what route (b) would be}

% ===========================================================================
\section{Two-dimensional results}
\label{sec:results2d}
% TRACE: notes §5.2 (#37), §5.3 (#39), §5.4 (#40), §5.5 (#41), §5.6 (#42);
% scripts/wave2d_stiff.py, scripts/wave2d_stiff_eigenvalues.py; numbers from
% #54's cache. Drafted by #59.
\subsection{Test problems and references}
\stub{\#59: 1-D spectral at normal incidence; Fourier-in-$x$ for oblique;
product-grid for the curved $\delta > 0$; the $122{,}500$-node jump-aware
run for the curved $\delta = 0$}
\subsection{Naive baseline}
\stub{\#59: the resolution floor hides the edge error in $v$; the spurious
$u$ separates unresolved from resolved edges $10\times$ more sharply}
\subsection{Stability}
\stub{\#59: spectra at $n = 900$ and $2500$, standard $\gamma$, every
$\delta$, flat and curved}
\subsection{The flat sweep and the rule}
\stub{\#59: seeds fourth order through $\delta = 0.0025$ at every $n$,
$20\times$ below naive in $v$ and $67\times$ in spurious $u$ at $19{,}600$
nodes; crossover at $h \approx \delta$; a resolved edge $1.4$--$2\times$
worse seeded; the rule: seed when $\delta \le h$}
\subsection{Oblique incidence}
\stub{\#59: seeds $1.4$--$1.9\times$ over naive at $26.6^\circ$, the
ablation worse than naive, every scheme about $2.5$ from the converted-S
floor, reported as it is}
\subsection{The curved edge}
\stub{\#59: jump stencils at the resolution floor at every $n$; seeds equal
the flat seeds at their own floor, $2.2$--$7\times$ ($\delta = 0.005$) and
$3$--$12\times$ ($\delta = 0.0025$) below naive; where route (a)'s geometry
first shows; the fine-end rate and the trim; costs}

% ===========================================================================
\section{Discussion}
\label{sec:discussion}
% TRACE: #27's comments (2026-09-19, 2026-09-20), notes §5.3 (the true
% Delta^3 of the seeds), §5.6.1 (seeded width), §5.6.2 (route (b), the
% curvilinear evaluation). Drafted by #60.
\subsection{Limitations}
\stub{\#60: FD4 / degree 3 only, one contrast, at most $19{,}600$ nodes, the
flat translation symmetry unexploited, one march per stencil, the port's
exclusions}
\subsection{Future work}
\stub{\#60: curvilinear evaluation of the marched seeds; seeded-width tuning;
the seeds' true $\Delta^3$ as hyperviscosity right-hand sides; M\"uhlbach's
recurrence in place of the dense solve; higher order; other contrasts; thin
bands; variable Lam\'e parameters; 3-D; the oblique train on the curved
geometry}
\subsection{Conclusions}
\stub{\#60: the rule and the one-construction claim, in \#53's wording}

% ===========================================================================
\section*{Acknowledgments}
% Decisions of 2026-09-20 (#51, recorded in #62): sole human author;
% acknowledgment to Bengt Fornberg for the 2016-era work this extends; explicit
% credit to Claude (Fable 5.1) in the form arXiv's policy on generative-AI
% tools asks for, checked at packaging time. #62 fixes the wording.
\stub{\#62: Bengt Fornberg, for the 2016-era work this construction extends
(the dissertation and the interface treatment of \cite{MartinFornberg2017});
the AI-assistance statement crediting Claude (Fable 5.1) for the
re-derivation, the code and the drafting}

% Skeleton only: render every seeded entry so the bibliography is known to
% compile. #53 verifies the entries; #61 removes \nocite{*}.
\nocite{*}
\bibliographystyle{amsalpha}
\bibliography{references}

\end{document}
